\chapter{Workflows}
\label{ch:workflows}

Workflows are the central abstraction in \haira{}. A workflow is a composable function with a trigger, typed parameters, and a typed return value. Workflows replace both n8n-style visual automations and LangGraph-style agent graphs.

\section{Workflow Declaration}

A workflow is a function preceded by a trigger decorator:

\begin{lstlisting}
@trigger_type("config")
workflow WorkflowName(param: Type, param: Type) -> ReturnType {
    // body
}
\end{lstlisting}

Workflows without a trigger decorator can be called as sub-workflows but cannot be triggered externally.

\section{Triggers}

Triggers define how a workflow is invoked. They use the \code{@} decorator syntax.

\subsection{@webhook --- HTTP Endpoint}

\begin{lstlisting}
@webhook("/summarize")
workflow Summarizer(url: string) -> Summary {
    content = fetch_page(url)
    summary = Researcher.ask("Summarize: {content}")
    return Summary{ text: summary, word_count: summary.len() }
}
\end{lstlisting}

The workflow parameters become the JSON request body. The return type becomes the JSON response body.

\begin{lstlisting}[language=bash,numbers=none]
$ curl -X POST http://localhost:8080/summarize \
    -d '{"url": "https://example.com/article"}'

{"text": "...", "word_count": 150}
\end{lstlisting}

Options:

\begin{lstlisting}
@webhook("/path")                  // POST (default)
@webhook("/path", method: "GET")   // GET
@webhook("/path", method: "PUT")   // PUT
\end{lstlisting}

\subsection{@websocket --- Bidirectional Streaming}

\begin{lstlisting}
@websocket("/chat")
workflow Chat(message: string, session_id: string) -> stream<string> {
    return Assistant.stream(message, session: session_id)
}
\end{lstlisting}

WebSocket workflows receive messages and can return \code{stream<string>} for real-time streaming responses.

\subsection{@cron --- Scheduled Execution}

\begin{lstlisting}
@cron("0 9 * * *")  // every day at 9 AM
workflow DailyReport() -> Report {
    data = fetch_metrics()
    report: Report, _ = Analyst.run(AnalyzeRequest{ data: data })
    send_email("team@company.com", "Daily Report", report.summary)
    return report
}
\end{lstlisting}

Cron expressions follow standard cron syntax (minute, hour, day-of-month, month, day-of-week).

\subsection{@event --- Event-Driven}

\begin{lstlisting}
@event("order.created")
workflow ProcessOrder(order: Order) -> OrderResult {
    validated = validate_order(order)
    if validated.has_errors {
        return OrderResult{ status: "rejected" }
    }
    charge_payment(order)
    return OrderResult{ status: "confirmed" }
}
\end{lstlisting}

Events are published using \code{haira.emit()}:

\begin{lstlisting}
haira.emit("order.created", Order{ id: "123", total: 99.99 })
\end{lstlisting}

\subsection{@manual --- CLI Only}

\begin{lstlisting}
@manual
workflow Migrate(version: string) -> { status: string } {
    // only invoked from CLI or main()
    run_migrations(version)
    return { status: "done" }
}
\end{lstlisting}

\begin{lstlisting}[language=bash,numbers=none]
$ haira run Migrate --input '{"version": "v2"}'
\end{lstlisting}

\subsection{No Trigger (Sub-Workflow)}

Workflows without a trigger can only be called from other workflows or from \code{main()}:

\begin{lstlisting}
workflow EnrichLead(email: string) -> EnrichedLead {
    company = lookup_company(email)
    social = lookup_social(email)
    return EnrichedLead{ email: email, company: company, social: social }
}
\end{lstlisting}

\section{Sequential Steps}

The workflow body is sequential by default. Each line can use the results of previous lines:

\begin{lstlisting}
@manual
workflow ArticlePipeline(topic: string) -> Article {
    // Step 1: Research
    research: ResearchResult, err = Researcher.run(ResearchRequest{ topic: topic })
    if err != nil { return Article{}, err }

    // Step 2: Write draft
    draft: Article, err = Writer.run(WriteRequest{
        content: research.summary,
        sources: research.sources
    })
    if err != nil { return Article{}, err }

    return draft
}
\end{lstlisting}

\section{Parallel Execution}

Use \keyword{spawn} blocks to run steps in parallel:

\begin{lstlisting}
@manual
workflow ReviewArticle(article: Article) -> [Review] {
    reviews = spawn {
        Reviewer.run(ReviewRequest{ article: article, focus: "accuracy" })
        Reviewer.run(ReviewRequest{ article: article, focus: "clarity" })
        Reviewer.run(ReviewRequest{ article: article, focus: "engagement" })
    }
    return reviews
}
\end{lstlisting}

A \keyword{spawn} block runs all statements concurrently and returns a list of results when all complete. This is equivalent to \code{Promise.all()} in JavaScript.

\section{Branching}

Standard \keyword{if}/\keyword{else} and \keyword{match} work inside workflows:

\begin{lstlisting}
@event("order.created")
workflow ProcessOrder(order: Order) -> OrderResult {
    validated = validate_order(order)

    if validated.has_errors {
        send_slack("#alerts", "Order {order.id} failed")
        return OrderResult{ status: "rejected" }
    }

    stock = check_stock(validated.items)

    if stock.available {
        charge_payment(order)
        send_email(order.customer_email, "Order confirmed!")
        return OrderResult{ status: "confirmed" }
    } else {
        send_email(order.customer_email, "Items out of stock")
        return OrderResult{ status: "backordered" }
    }
}
\end{lstlisting}

\section{Loops}

\begin{lstlisting}
@manual
workflow ProcessBatch(urls: [string]) -> [string] {
    results = for url in urls {
        content = fetch_page(url)
        Researcher.ask("Summarize: {content}")
    }
    return results
}
\end{lstlisting}

Loops inside workflows are sequential by default. For parallel iteration, use \keyword{spawn}:

\begin{lstlisting}
@manual
workflow ProcessBatchParallel(urls: [string]) -> [string] {
    results = spawn {
        for url in urls {
            fetch_page(url) | Researcher.ask("Summarize: {_}")
        }
    }
    return results
}
\end{lstlisting}

\section{Sub-Workflows}

Workflows can call other workflows like functions:

\begin{lstlisting}
workflow ScoreLead(company: Company, social: SocialProfile) -> ScoreResult {
    score: ScoreResult, err = Analyst.run(ScoreRequest{
        company: company,
        social: social
    })
    return score
}

workflow EnrichLead(email: string) -> EnrichedLead {
    company = lookup_company(email)
    social = lookup_social(email)
    score = ScoreLead(company, social)
    return EnrichedLead{ email: email, company: company, score: score.value }
}

@webhook("/new-lead")
workflow SalesAutomation(email: string) -> { result: string } {
    lead = EnrichLead(email)

    if lead.score > 0.8 {
        send_slack("#sales", "Hot lead: {email}")
    }

    return { result: "processed" }
}
\end{lstlisting}

Sub-workflows compose naturally. The compiler validates type compatibility at each call site.

\section{Error Handling}

Workflows use standard \haira{} error handling:

\begin{lstlisting}
@manual
workflow SafeProcess(data: string) -> { result: string } {
    result, err = risky_operation(data)
    if err != nil {
        send_slack("#errors", "Failed: {err.message}")
        result = fallback_operation(data)
    }
    return { result: result }
}
\end{lstlisting}

\section{Running Workflows}

\subsection{From main()}

\begin{lstlisting}
fn main() {
    result, err = ArticlePipeline("AI agents")
    if err != nil {
        io.println("Failed: " + err.message)
        return
    }
    io.println(result.title)
}
\end{lstlisting}

\subsection{Server Mode}

Start a server that listens for all triggered workflows:

\begin{lstlisting}
fn main() {
    server = http.Server([
        Summarizer,         // POST /summarize
        Chat,               // WebSocket /chat
        ProcessOrder,       // on order.created events
        DailyReport,        // cron 0 9 * * *
    ])

    io.println("Haira server running on :8080")
    server.listen(8080)
}
\end{lstlisting}

\subsection{CLI}

\begin{lstlisting}[language=bash,numbers=none]
# Run a workflow directly
$ haira run ArticlePipeline --input '{"topic": "AI agents"}'

# Start server mode
$ haira serve main.haira --port 8080

# Run with environment
$ ANTHROPIC_API_KEY=sk-... haira serve main.haira
\end{lstlisting}

\section{Workflow Grammar}

\begin{lstlisting}[language=EBNF]
workflow_decl = { decorator } "workflow" identifier
                "(" [ params ] ")" [ "->" type ] block ;

decorator     = "@" identifier [ "(" [ arguments ] ")" ] ;

spawn_block   = "spawn" "{" { statement } "}" ;
\end{lstlisting}
